\chapter{Complexity — Meyer Tab}
%%% generated by the blueprint pipeline from TCSlib.Complexity.CircuitComplexity.MeyerTab
\section{Overview}

Declarations of the Lean module \texttt{TCSlib.Complexity.CircuitComplexity.MeyerTab}.

\section{Declarations}

\begin{definition}[Answer from an iterated configuration]
\lean{Complexity.Meyer.tabAns}
\label{Complexity.Meyer.tabAns}
\leanok
\uses{Complexity.Meyer.Sel, Complexity.Meyer.selAnswer, Complexity.Meyer.wstep, Complexity.TimeHierarchy.OutReg.ofList, Turing.Cfg, Turing.Cfg.inputSymbol, Turing.Cfg.workTapeSymbols}
Let $M$ be a finite Turing machine with Boolean symbols, let a selector and a
configuration relative to a Boolean input be given, and let $a\in\mathbb{N}$. The
Boolean value is the selector’s answer on the configuration’s state, its output
register, and the input and work-tape symbols in the configuration obtained after
applying the selected work step $a$ times.
\end{definition}

\begin{definition}[Boolean function from a machine encoding]
\lean{Complexity.Meyer.tabFun}
\label{Complexity.Meyer.tabFun}
\leanok
\uses{Complexity.Meyer.decodeSel, Complexity.Meyer.nbits, Complexity.Meyer.parseQ, Complexity.Meyer.tabAns, Complexity.TimeHierarchy.ctrVal, Turing.MultiTapeTM.initCfg, Turing.MultiTapeTM.runFrom}
Let $M$ be a finite Turing machine with Boolean symbols and let $y$ be a finite Boolean
string. If $y$ parses as an encoding of selector bits, a time parameter, an iteration
parameter, and an input $x$, return the answer obtained by decoding the selector bits
and evaluating the machine from its initial configuration on $x$ for the encoded time,
then taking the selector’s answer after the encoded number of work steps. If $y$ does
not parse in this way, return false.
\end{definition}

\begin{definition}[Language accepted by the table function]
\lean{Complexity.Meyer.Tab}
\label{Complexity.Meyer.Tab}
\leanok
\uses{Complexity.Meyer.tabFun}
Let $M$ be a finite Turing machine with Boolean symbols. The language consists of
exactly those finite Boolean strings on which the Boolean function defined from $M$’s
machine encoding returns true.
\end{definition}

\begin{theorem}[Bound on the Value of a Boolean Word]
\lean{Complexity.Meyer.ctrVal_lt}
\label{Complexity.Meyer.ctrVal_lt}
\leanok
\uses{Complexity.TimeHierarchy.ctrVal}
\difficulty{4}
For every finite Boolean word $w$, its binary value is strictly less than $2^{|w|}$.
\end{theorem}

\begin{theorem}[Lengths of Parsed Fields]
\lean{Complexity.Meyer.parseQ_length}
\label{Complexity.Meyer.parseQ_length}
\leanok
\uses{Complexity.Meyer.nbits, Complexity.Meyer.parseQ, Complexity.Meyer.readField, Complexity.Meyer.readField_eq_some}
\difficulty{5}
Let $M$ be a finite Turing machine over the Boolean alphabet, and let $y$ be a Boolean
word. If parsing $y$ with the number of bits determined by the state set and number of
work tapes of $M$ succeeds, yielding a bit vector and three words $t,a,x$, then $|t|\le
|y|$ and $|a|\le |y|$.
\end{theorem}

\begin{theorem}[Initial input-head boundary tag]
\lean{Complexity.Meyer.init_tag}
\label{Complexity.Meyer.init_tag}
\leanok
\uses{Turing.FinTM.VirtualTag, Turing.MultiTapeTM.initCfg}
\difficulty{1}
Let $M$ be a finite Turing machine over the Boolean alphabet, and let $x$ be a finite
Boolean input. In the initial configuration of $M$ on $x$, the input head is at position
$1$ and its Boolean boundary tag is true: the tag is false at position $0$ and true at
the right boundary position $|x|+1$.
\end{theorem}

\begin{theorem}[Tabulation Machine Computes the Function]
\lean{Complexity.Meyer.tabTM_computes}
\label{Complexity.Meyer.tabTM_computes}
\leanok
\uses{Complexity.Meyer.acfg_eq_lcfg, Complexity.Meyer.ctrVal_lt, Complexity.Meyer.fpos, Complexity.Meyer.init_tag, Complexity.Meyer.lcfg, Complexity.Meyer.loop_run, Complexity.Meyer.nbits, Complexity.Meyer.parseQ, Complexity.Meyer.parseQ_length, Complexity.Meyer.tabAns, Complexity.Meyer.tabFun, Complexity.Meyer.tabTM, Complexity.Meyer.tabTM_setup, Complexity.Meyer.walk_run, Complexity.TimeHierarchy.OutReg.ofList, Complexity.TimeHierarchy.ctrVal, Turing.Cfg.init, Turing.FinTM.ComputesFunInTime, Turing.FinTM.ComputesInTime.mono, Turing.FinTM.bufferTape, Turing.FinTM.computesInTime_iff, Turing.MultiTapeTM.initCfg, Turing.MultiTapeTM.runFrom, Turing.MultiTapeTM.runFrom_add}
\difficulty{6}
For every Boolean finite Turing machine $M$, the tabulation machine built from $M$
computes the function that sends each Boolean word $y$ to the one-symbol output
consisting of the tabulation function’s value on $y$, within time bound $10(n+1)2^n$ on
inputs of length $n$.
\end{theorem}

\begin{theorem}[Exponential-time membership of the tabulation language]
\lean{Complexity.Meyer.Tab_mem_EXP}
\label{Complexity.Meyer.Tab_mem_EXP}
\leanok
\uses{Complexity.EXP, Complexity.Meyer.Tab, Complexity.Meyer.tabFun, Complexity.Meyer.tabTM, Complexity.Meyer.tabTM_computes, Turing.FinTM.ComputesInTime.mono, Turing.MultiTapeTM.indicator}
\difficulty{4}
For every finite-state Turing machine $M$ over the Boolean alphabet, the language of
Boolean words accepted by the tabulation procedure associated with $M$ belongs to
exponential time.
\end{theorem}
